\documentclass[10pt]{article}
\usepackage{multicol}
\usepackage{fontawesome5}
\usepackage{xltxtra}
\usepackage{marathi}
\usepackage{kantlipsum}
\टंक{Eczar}
\usepackage{xcolor}
\usepackage{hyperref}
\hypersetup{%
  colorlinks,
  linkcolor={red!50!black},
  citecolor={blue!50!black},
  urlcolor={blue!80!black}
}%
\title{एग्झार}
\author{}
\date{%
  आवृत्ती ०.१ - १० फेब्रुवारी, २०२१
  \hfill
  संस्करण ०.१ - १० फ़रवरी, २०२१\\
  \bigskip
  {\small\faIcon{gitlab}\quad\url{https://github.com/rosettatype/eczar}}
}%
\sloppy
\hyphenpenalty=10000
\exhyphenpenalty=1000
\setlength\columnsep{20pt}

\begin{document}
\maketitle

\vfill

\अंतरबदल{1.5}
\begin{abstract}
  \begin{multicols}{2}
    रोझेटा संस्थेचा युनिकोड-आधारित एग्झार हा टंक टेक्-वितरणात उपलब्ध करून देण्याकरिता हा
    आज्ञासंच सादर करण्यात येत आहे. ह्या टंकाचा समावेश टेक्-वितरणात झाल्यास एग्झार हा टंक
    फॉन्टस्पेकसदृश आज्ञासंचांसह सहज वापरता येईल. ह्या दस्तऐवजात एग्झार टंकाचे दृश्य-उदाहरण सादर
    करण्यात आले आहे.
    
    \columnbreak
    रोज़ेटा संस्था द्वारा प्रस्तुत युनिकोड आधारित एग्ज़ार नामक टंक का लाटेक् में प्रयोग सुलभ हो इस
    लिए यह आज्ञासमुच्चय प्रस्तुत किया गया है। टेक् वितरण में समाविष्ट होने के पश्चात् फॉन्टस्पेक जैसे
    आज्ञासमुच्चयों के साथ इस टंक का प्रयोग लाटेक् में हो सकेगा। इस दस्तावेज़ में एग्ज़ार टंक का दृश्य
    उदाहरण प्रस्तुत किया है।
  \end{multicols}

  \paragraph{Abstract:} This package provides Eczar, a Unicode-based font family
  which supports $45+3$ languages with two scripts i.e.\ Devanagari and
  Latin. This font can be used with pakages like \verb|fontspec| \& \XeLaTeX\ or
  Lua\LaTeX.
\end{abstract}
\vfill
\clearpage\pagebreak

\begin{multicols}{2}
  \section*{नमुना पाठ्य मराठी}
  \परिच्छेद

  \columnbreak
  \section*{नमूना पाठ्य हिन्दी}
  नमस्ते! यह पाठ्य अर्थशून्य है। यहाँ क्या और कैसे मुद्रित होगा इस का यह सिर्फ एक नमूना है। इसे
  पढ़ने से आप को कोई जानकारी नहीं मिलेगी। सच? क्या इस में कोई जानकारी नहीं है? इस पाठ्य में और
  `पिढ्ढ करढपाखू' ऐसे निरर्थक शब्दों में क्या कोई अंतर है? हाँ! इसे अन्धपाठ्य कहते हैं। अन्धपाठ्य से
  चुना गया टंक कौन सा है, उस में अक्षर कैसे दिखते हैं, इन विषयों के बारे में जानकारी प्राप्त होती है।
  अन्धपाठ्य के लिए विशिष्ट शब्दों की आवश्यकता नहीं, किन्तु शब्द चुनी गयी भाषा के होना आवश्यक हैं।
\end{multicols}
\section*{Sample text English}
\kant[1]
\end{document}